%% ============================================================================
\section{Introduction}
\label{sec:intro}
%% ============================================================================

% Societal and economic importance of rolling element bearings and their lubrication
Rolling element bearings are among the most prevalent mechanical components in industrial machinery, and their reliable operation depends critically on adequate lubrication. Lubrication failure---whether caused by starvation, viscosity degradation, or contamination---increases metal-to-metal asperity contact, accelerates surface fatigue and wear, raises operating temperature, and ultimately leads to premature bearing failure. Improper lubrication is estimated to account for between 30\% and 80\% of all rolling element bearing failures~\cite{jakobsen_detecting_2023,dai_situ_2024,nassef_impact_2022}.

% Research gap in lubrication condition monitoring (LCM)
Despite this, lubrication condition monitoring (LCM) has received substantially less research attention than vibration-based fault detection and remaining useful life estimation, which are now mature technologies deployed widely in industry\todo{Reference for this}. The principal reason is that lubrication-related signal changes are subtler and more entangled with operating conditions than the impulse signatures of surface spalling or raceway cracks\todo{Reference for this}. Accurate, online LCM would enable condition-based relubrication: applying lubricant at the right time in the right quantity rather than following time-based schedules that systematically produce both under- and over-lubrication~\cite{wakiru_review_2019}.

% Conventional CM approaches and their limitations for LCM
The dominant approach to bearing condition monitoring is vibration analysis, which is highly effective at detecting macroscopic faults such as raceway spalling and rolling element damage through characteristic frequency analysis. However, vibration-based methods are inherently less sensitive to the subtle tribological changes associated with lubrication film formation and breakdown: under varying speed and load, operating-condition effects dominate the spectral signature and mask the comparatively weak lubrication-dependent component~\cite{shang_comprehensive_2024,cornel_condition_2021}. Temperature monitoring and oil analysis provide complementary information but are limited by indirect measurement, sampling delay, or the need for access to the lubricant reservoir\todo{Reference? Seems like a very specific mention which might stem from an actual article}. None of these approaches yields a direct, continuous, and non-invasive indicator of the instantaneous lubrication regime under operational conditions.

% Specifically acoustic emission and passive ultrasound
Passive, non-invasive sensing methods---acoustic emission (AE) and passive ultrasound (US)---are practically attractive for industrial LCM because they require no bearing modification and can be retrofitted to bearings in service. AE is sensitive to the broadband stress waves emitted by asperity contact events and microcrack propagation across a wide frequency range (typically 50~kHz--\SI{4}{MHz})\todo{Reference}. Passive ultrasound in the 38--40~kHz band is already industrially deployed in commercial instruments (UE Systems, Sonotec)\todo{Reference} for threshold-based lubrication assessment: the broadband ultrasonic level is compared against a reference baseline and relubrication is triggered when the level rises by a specified number of decibels~\cite{lineham_ultrasonic_2008,jakobsen_detecting_2023,jakobsen_low_2024}. Despite both sensing approaches being in practical use, their signal-processing behaviour has not been characterised simultaneously in a unified regression framework with a physics-grounded lubrication reference variable.

% Literature review of AE and passive US for LCM
The sensitivity of AE to lubrication regime has been established across several bearing configurations. Miettinen et al.~\cite{miettinen_analysis_2001} showed that AE pulse count and RMS distinguish boundary, mixed, and full-film lubrication regimes and are sensitive to base oil viscosity and thickener content in grease-lubricated bearings. Yoshioka and Shimizu~\cite{yoshioka_compound_2009} demonstrated a compound AE and vibration monitoring system on a grease-lubricated deep-groove ball bearing, using envelope demodulation and statistical AE parameters to characterise lubrication state. Tandon et al.~\cite{tandon_ae_2006} found that AE and shock-pulse sensors outperform accelerometers for detecting grease contamination in a ball bearing---the sensing scenario most closely related to the present study. More recently, Cornel et al.~\cite{cornel_condition_2021} showed that AE detects lubrication-sensitive surface damage such as micro-pitting and smearing up to eight hours earlier than vibration monitoring, and Hou and Zhang~\cite{hou_insitu_2025} derived amplitude-centre and energy-centre features that classify lubrication status across speeds and viscosity levels on a thrust rolling bearing, without a continuous physics-grounded regression target. On signal processing, Renhart et al.~\cite{renhart_ae_2021} proposed octave filter-bank decomposition with RMS--friction product features for AE-based tribological event characterisation in bearing contacts---applying frequency sub-band analysis to contact-type identification rather than to continuous lubrication adequacy regression---and Marticorena and Garc\'{i}a Peyrano~\cite{marticorena_cage_2022} showed that spectral correlation of AE signals resolves lubrication regime changes on a ball bearing under varying load, demonstrating that kinematic and acoustic features together can partially decouple operating-condition confounds. The physical link between lubrication film adequacy and the resulting AE signature has been formalised by Liu et al.~\cite{liu_combined_2021} and Zhang et al.~\cite{zhang_coupled_2026}, who developed coupled vibration--AE models incorporating EHL film stiffness and lubricant traction dynamics, establishing that emission intensity and frequency content are governed by the film thickness and the rate of asperity contact. For the passive ultrasound channel, Jakobsen et al.~\cite{jakobsen_detecting_2023,jakobsen_low_2024} demonstrated that broadband acoustic features from a MEMS microphone sensor carry measurable information about the ISO~281 viscosity ratio $\kappa$ and established the $\kappa$-regression framework that this paper extends. Despite this body of work, no study has simultaneously acquired calibrated AE and passive US from a single ball bearing and quantitatively assessed their complementarity for $\kappa$ regression under a common experimental protocol---the gap this paper addresses.

% Regression to kappa in literature
A key limitation of existing data-driven LCM work is the reliance on discrete fault labels or qualitative lubrication states rather than physically interpretable, continuous targets. The ISO~281 viscosity ratio $\kappa$ addresses this directly: it encodes the ratio of the actual lubricant kinematic viscosity at operating temperature to the minimum viscosity required for full elastohydrodynamic (EHL) film formation at the current speed and bearing geometry, and thereby maps the instantaneous operating conditions to a standardised lubrication adequacy index. A model trained to predict $\kappa$ implicitly learns the full multi-dimensional dependence on speed, temperature, and lubricant properties, and its output is directly interpretable in terms of established lubrication regime boundaries. Jakobsen et al.~\cite{jakobsen_detecting_2023,jakobsen_low_2024} demonstrate $\kappa$ regression from vibration and passive US features using LASSO and neural networks, establishing the methodological foundation that this paper builds upon. Despite its relevance as a continuous, physics-grounded target, the use of $\kappa$ as a regression objective in acoustic condition monitoring remains largely unexplored.

% Specific research gap for AE and US fusion
The complementarity of AE and passive US for LCM has been argued on physical grounds in reviews~\cite{shang_comprehensive_2024,wang_advanced_2024}, but has not been quantified experimentally in a unified framework. Whether jointly exploiting both sensor modalities improves $\kappa$ prediction over either alone is an open question with direct practical implications: additional sensors impose cost and complexity, and are only justified if they provide non-redundant information.

% Scope statement
The present study is conducted under controlled conditions using a single bearing configuration, lubricant, and loading, with operating variations limited to rotational speed and temperature. This scope is a deliberate choice: controlled conditions allow lubrication-dependent acoustic behaviour to be isolated from confounds arising from bearing-to-bearing variability and degradation progression. Generalisability to other bearing types, lubricants, and operating environments is an acknowledged limitation and a primary direction for future work.

\subsection*{Contributions}

This paper makes the following contributions:
\begin{enumerate}
  \item \textbf{Novel dual-channel experiment.} The first study to simultaneously acquire AE and passive US signals from a ball bearing and regress both, individually and jointly, to $\kappa$ using a common dataset and analytical ground truth.
  \item \textbf{Sub-band feature extraction for AE.} Features are computed from the broadband AE signal and from three frequency sub-bands (20--500~kHz, 500~kHz--1~MHz, 1--2~MHz), enabling model-driven identification of the most informative AE frequency range for lubrication state prediction.
  \item \textbf{Quantified fusion test.} The first quantitative test of AE--US complementarity in a unified experimental framework, including a positive fusion result that directly addresses an open question in the LCM literature.
  \item \textbf{Rigorous comparative evaluation.} All model comparisons use repeated nested cross-validation with corrected significance tests (Nadeau--Bengio, Holm--Bonferroni, Wilcoxon, Diebold--Mariano, bootstrap $\Delta$RMSE) that account for the correlation structure of cross-validated estimates, converting the fusion and sub-band comparisons into statistically grounded conclusions rather than point estimates susceptible to inflated confidence.
\end{enumerate}

The remainder of the paper is organised as follows. Section~\ref{sec:background} reviews the theoretical background of $\kappa$ and the physical origins of AE and UL signals. Section~\ref{sec:setup} describes the experimental setup and data collection. Section~\ref{sec:pipeline} presents the signal processing and feature extraction pipeline. Section~\ref{sec:features} presents the feature analysis results. Section~\ref{sec:modelling} presents the regression modelling and statistical evaluation. Section~\ref{sec:discussion} discusses the findings. Section~\ref{sec:conclusion} concludes.
